\documentclass[12pt]{article}
\usepackage[T1]{fontenc}
\usepackage[utf8]{inputenc}
\usepackage{lmodern}
\usepackage{amsmath,amssymb}
\usepackage{graphicx}
\usepackage{hyperref}
\usepackage[a4paper,margin=20mm]{geometry}
\setlength{\parindent}{0pt}
\setlength{\parskip}{0.25em}
\begin{document}
\begin{center}{\LARGE\textbf{Physics problem set}}\end{center}

\begin{enumerate}
\setlength{\itemsep}{0.5em}
\setlength{\topsep}{0.4em}
\setlength{\parsep}{0.2em}
\item Egy kerékpáros „teljes erőbedobással” lejtőn felfelé $v_{1}=12 \mathrm{~km} / \mathrm{h}$, ugyanezen lejtőn lefelé $v_{2}=36 \mathrm{~km} / \mathrm{h}$ sebességgel tud haladni. Mekkora a kerékpáros legnagyobb sebessége vízszintes úton, ha a maximális erőkifejtése független a sebességétől?


A kicsit pongyolán megfogalmazott „erőkifejtés” szó egy fizikus értelmezésében jelentheti
a) a kerékpáros által kifejtett (és azt a hajtókarok, a lánckerekek és a lánc által a kerekekhez továbbított) erő nagyságát;
b) a kerékpáros mechanikai teljesítményét.


Gondoljuk végig a feladatot mindkét értelmezésben!
\hfill{\small\textit{Source: 333+ Physics Problems · 2017 · F. 38.}}
\item Andi és Bandi „gurulásos versenyt” rendeznek egy elegendően hosszú, állandó hajlásszögü lejtőn. Mindkettőjük biciklije ugyanolyan, és egyikük sem hajtja a kerékpárját. Andi 50 kg, Bandi 100 kg tömegú, a kerékpár tömege pedig 10 kg. Bandi légellenállása (a túlsúlyával összefüggő nagyobb „mérete“ miatt) másfélszer akkora, mint - ugyanakkora sebesség esetén - Andié. Melyikük gurul messzebbre a lejtőhöz csatlakozó vízszintes úton?
\begin{center}
\includegraphics[width=0.7\linewidth]{\detokenize{figures/image-001.png}}
\end{center}
(Feltételezhetjük, hogy a kerékpárosokat a sebesség négyzetével arányos közegellenállási erő, továbbá a csapágyak súrlódása és a gördülő ellenállás fékezi. Az utóbbiakat kezelhetjük úgy, mintha valamekkora $\mu$ „effektív” súrlódási együtthatónak megfelelő csúszó súrlódás lépne fel.)
\hfill{\small\textit{Source: 333+ Physics Problems · 2017 · F. 39.}}
\item Vízszintes asztallapon egy terítő, rajta pedig egy acélgolyó található. A terítőt kihúzzuk a golyó alól, eközben a golyó a súrlódási erón miatt mozgásba jön, és forogni is kezd. Mekkora lesz a sebessége az asztalon, amikor már csúszásmentesen gördül? (Az asztal elég nagy méretú, a golyó nem esik le róla.)
\begin{center}
\includegraphics[width=0.7\linewidth]{\detokenize{figures/image-002.png}}
\end{center}

\begin{center}
\includegraphics[width=0.7\linewidth]{\detokenize{figures/image-003.png}}
\end{center}
\hfill{\small\textit{Source: 333+ Physics Problems · 2017 · F. 99.}}
\end{enumerate}
\end{document}
